\documentclass{article}
\usepackage{amsmath,amssymb}
\usepackage{xurl}
\usepackage[colorlinks=true,linkcolor=blue,citecolor=blue,urlcolor=blue]{hyperref}
% Shared commands for notes in docs/paper-gaps/.

\DeclareUnicodeCharacter{2080}{\ifmmode{}_{0}\else\textsubscript{0}\fi}
\DeclareUnicodeCharacter{2081}{\ifmmode{}_{1}\else\textsubscript{1}\fi}
\DeclareUnicodeCharacter{2082}{\ifmmode{}_{2}\else\textsubscript{2}\fi}
\DeclareUnicodeCharacter{2099}{\ifmmode{}_{n}\else\textsubscript{n}\fi}

\newcommand{\Lean}{\textsc{Lean}}
\ExplSyntaxOn
\NewDocumentCommand{\leanid}{m}{
  \ifmmode
    \textsf{\detokenize{#1}}
  \else
    \group_begin:
    \urlstyle{sf}
    \tl_set:Nn \l_tmpa_tl {#1}
    \tl_replace_all:Nnn \l_tmpa_tl {\_} {_}
    \exp_args:NV \path \l_tmpa_tl
    \group_end:
  \fi
}
\ExplSyntaxOff
\providecommand{\tr}{\operatorname{tr}}
\newcommand{\tnleanIssueBase}{https://github.com/LionSR/TNLean/issues/}
\newcommand{\tnleanPrBase}{https://github.com/LionSR/TNLean/pull/}
\newcommand{\tnleanDocBase}{https://sirui-lu.com/TNLean/docs/}
\newcommand{\ghissue}[1]{\href{\tnleanIssueBase#1}{GitHub issue \##1}}
\newcommand{\ghpr}[1]{\href{\tnleanPrBase#1}{PR \##1}}
\newcommand{\doclink}[1]{\href{\tnleanDocBase#1.html}{\path{#1}}}

% Dirac notation and the identity operator, as in blueprint/src/macros/common.tex.
% \providecommand keeps note-local definitions authoritative.
\usepackage{dsfont}
\providecommand{\ket}[1]{|#1\rangle}
\providecommand{\bra}[1]{\langle #1|}
\providecommand{\braket}[2]{\langle #1 | #2 \rangle}
\providecommand{\ketbra}[2]{|#1\rangle\!\langle #2|}
\providecommand{\Id}{\mathds{1}}
\providecommand{\one}{\mathds{1}}
\providecommand{\C}{\mathbb{C}}
\providecommand{\MN}[1]{M_{#1}(\C)}
\providecommand{\MD}{M_D(\C)}

% Verdict marker: \gapnote{<kind>}{<status>}. Kind is the policy
% classification (clarification, local-correction, scope-restriction,
% unfaithful, false-source, open-gap); status is open, wip (elimination
% underway), resolved, or historical. Machine-read by the site index;
% prints a verdict line.
\newcommand{\gapnote}[2]{\par\noindent\textbf{Verdict:} #1 (#2).\par}

\providecommand{\leanid}[1]{\textsf{\detokenize{#1}}}
\providecommand{\doclink}[1]{%
    \href{https://sirui-lu.com/TNLean/docs/#1.html}{\path{#1}}}
\providecommand{\gapnote}[2]{%
    \par\noindent\textbf{Verdict:} #1 (#2).\par}
\title{Numerical Scanner Conditions at the Fixed Bootstrap Scales}
\author{}
\date{October 14, 2026}
\begin{document}
\maketitle
\gapnote{scope-restriction}{open}

For $0<e_0<1$, the source fixes
$g_0=(1-e_0)/2$, $\mu=1-g_0$, $\nu=g_0/1000$,
$\ell=\nu/100$, $\kappa=\min\{\ell,e_*/4\}$ and $W=4$ once
for the whole bootstrap.
They satisfy the scanner exponent conditions, and their rounded scales
remain
\begin{align}
    L&=\lfloor n^{1-\ell}\rfloor,
    &m&=\lfloor n^\mu\rfloor,
    &D&=\lceil n^\kappa\rceil,\nonumber\\
    K&=\left\lfloor L/(8m)\right\rfloor,
    &a&=4/K,
    &r_0&=\lceil C_r\log^2n\rceil.
    \label{eq:bootstrap_numerics_scales}
\end{align}
The numerical datum assigns its auxiliary entropy exponent the value
$\mu$; this assignment does not constrain the actual entropy exponent
being improved.
\footnote{Formal definition:
    \leanid{TNLean.PEPS.AreaLaw.Scan.BootstrapParameters.scannerExponents}
    in \doclink{TNLean/PEPS/AreaLaw/Scan/BootstrapScannerExponents}.}

For every scanner exponent tuple and every fixed $q\geq1$, $C_r\geq0$
and $c_t>0$, one threshold gives $K\geq8$, $a>0$ and
$a\ell_q(r_0)\leq c_t$.
Indeed, $\alpha=1-\ell-\mu>0$, the literal band count satisfies
$n^\alpha\leq C_KK$, and the actual physical logarithmic dimension
satisfies $\ell_q(r_0)\leq C_d\log^4n$.
Thus $a\ell_q(r_0)\leq4C_KC_d\log^4n/n^\alpha\to0$.
No marginal moment or sharp-comparison margin is used.
\footnote{Formal declaration:
    \leanid{TNLean.PEPS.AreaLaw.Scan.ScannerExponents.exists_fixedWeightFour_transport_smallness}
    in \doclink{TNLean/PEPS/AreaLaw/Scan/FixedWeightTransportSmallness}.}

A further generic threshold gives $2r_0\leq D$, $1\leq D$,
$4D\leq m$, $8\leq K$ and $8Km\leq L$ with precisely the same
rounded scales.
The first estimate includes the doubled radius used by the physical
truncation.
The existing rounded-clearance estimate applies to these exponents as
well: for fixed positive $C_2$ and $n=\lceil C_2s\rceil$, one
threshold in the parent side length precedes every positive safety
parameter and gives the original safe-box clearance budgets.
\footnote{Formal declaration:
    \leanid{TNLean.PEPS.AreaLaw.Scan.ScannerExponents.exists_rounded_admissibility}
    in \doclink{TNLean/PEPS/AreaLaw/Scan/ScannerRoundedAdmissibility}.
    The existing clearance theorem is
    \leanid{TNLean.PEPS.AreaLaw.Scan.ScannerExponents.exists_rounded_clearance_threshold}
    in \doclink{TNLean/PEPS/AreaLaw/Scan/RoundedClearance}.}

Every threshold is chosen after its fixed scalar coefficients and
exponents and before the system size, finite physical instance or
actual scanner data.
No uniform threshold over arbitrarily small positive exponent gaps is
claimed.
These results supply numerical prerequisites only.
They do not construct a scanner or physical state, prove entropy
transport, select an actual interpolation, or carry out the remaining
bootstrap entropy absorption.
The source's final-stage exponents are not substituted into the fixed
bootstrap parameters.
No full scanner theorem, Proposition~9.6 or area law is asserted.
\footnote{Source: Ref.~\cite{OpenAI2026AreaLaw},
    \texttt{08-scanner.tex}, lines 32--49, 358--359 and 693--755,
    revision \texttt{adc7f1241b42e322a6451854ab7e4b4c146bf78a}.}

\bibliographystyle{alpha}
\bibliography{references}
\end{document}
